% --- Archon physics formalization source begin ---
% archon:physics
% archon:covers PhyXMiniProblems/problem_phyx_mini_0810.lean
% archon:source-report reports/phyx_mini/problem_phyx_mini_0810.source.json

\chapter{Physics problem phyx\_mini\_0810}
\label{ch:PhyXMiniProblems_problem_phyx_mini_0810}

\paragraph{Problem source.}
\#\# Physical scenario

You throw a 0.145 kg baseball straight up, giving it an initial velocity of magnitude 20.0 m/s.

\#\# Question

Find how high it goes, ignoring air resistance.

\#\# Answer choices

- A: A: 23.1m

- B: B: 20.4m

- C: C: 27.2m

- D: D: 16.5m

\#\# Figure caption (auxiliary; use the image as primary evidence)

The image illustrates a physics problem involving a baseball being thrown vertically. 

- At the bottom, a hand is shown throwing a baseball upwards with a velocity \textbackslash{}( v\_1 = 20.0 \textbackslash{}, \textbackslash{}text\{m/s\} \textbackslash{}).
- The baseball has a mass \textbackslash{}( m = 0.145 \textbackslash{}, \textbackslash{}text\{kg\} \textbackslash{}).
- The initial position of the baseball is marked as \textbackslash{}( y\_1 = 0 \textbackslash{}).
- Toward the top, the baseball reaches a point where its velocity \textbackslash{}( v\_2 = 0 \textbackslash{}), which represents the peak of its trajectory.
- This peak position is labeled as \textbackslash{}( y\_2 \textbackslash{}). 
- An arrow pointing upwards from the hand represents the direction of the throw.
- The baseball is depicted twice: once at the hand and once at the peak.

The scene models the motion of the baseball under gravitational effects from an initial upward throw to its highest point.

\#\# Dataset metadata

Kinematics; Physical Model Grounding Reasoning, Multi-Formula Reasoning

\paragraph{Recorded answer/context.}
B: B: 20.4m

\paragraph{Figure/image path.}
/root/proposal\_for\_physic/science-mango/phyx\_mini\_run/phyx\_data/test\_image/810.png

\paragraph{Formalization target.}
create a compiling Lean file with sorry bodies at `PhyXMiniProblems/problem\_phyx\_mini\_0810.lean`. The Lean declarations must preserve the physical quantities, dimensions or dimensional roles, figure labels, governing-law hypotheses, and final relation expressed by this problem.
Use Mathlib/Physlib names found through LeanExplore where available. If a physics API is missing, introduce faithful local abstractions rather than scalar placeholder aliases.

\begin{theorem}[Physics formalization target]
\label{thm:physics:phyx_mini_0810:target}
\lean{PhyXMiniProblems.ProblemPhyXMini0810.maximumHeight_matches_answerB}
\uses{def:physics:phyx-mini-0810:phyxminiproblems-problemphyxmini0810-baseballverticalthrowsetup, def:physics:phyx-mini-0810:phyxminiproblems-problemphyxmini0810-heightgaininmeters, def:physics:phyx-mini-0810:phyxminiproblems-problemphyxmini0810-matchesproblemdescription, def:physics:phyx-mini-0810:phyxminiproblems-problemphyxmini0810-matchesprimaryfigure, def:physics:phyx-mini-0810:phyxminiproblems-problemphyxmini0810-usesstandardearthgravity, def:physics:phyx-mini-0810:phyxminiproblems-problemphyxmini0810-hasphysicalthrowparameters, def:physics:phyx-mini-0810:phyxminiproblems-problemphyxmini0810-satisfiesuniformgravitykinematics, def:physics:phyx-mini-0810:phyxminiproblems-problemphyxmini0810-fitsdisplayedheight}
The assigned autoformalize agent should translate this physics problem into Lean declarations in the covered file, with theorem and lemma proof bodies written as `by sorry`.
\end{theorem}
\begin{proof}
This is an autoformalization task, not a proof task. Produce statements that can later be proved by the physics prover without weakening the physical model.
\end{proof}

% --- Archon declaration topology begin ---
\section{Declaration topology}
\begin{definition}[vertical Velocity Dimension]
  \label{def:physics:phyx-mini-0810:phyxminiproblems-problemphyxmini0810-verticalvelocitydimension}
  \lean{PhyXMiniProblems.ProblemPhyXMini0810.verticalVelocityDimension}
  The physical dimension of signed vertical velocity, L T⁻¹
\end{definition}
\begin{proof}
  This is the declared model definition of vertical Velocity Dimension.
\end{proof}
\begin{definition}[vertical Acceleration Dimension]
  \label{def:physics:phyx-mini-0810:phyxminiproblems-problemphyxmini0810-verticalaccelerationdimension}
  \lean{PhyXMiniProblems.ProblemPhyXMini0810.verticalAccelerationDimension}
  The physical dimension of signed vertical acceleration, L T⁻²
\end{definition}
\begin{proof}
  This is the declared model definition of vertical Acceleration Dimension.
\end{proof}
\begin{definition}[Mass Quantity]
  \label{def:physics:phyx-mini-0810:phyxminiproblems-problemphyxmini0810-massquantity}
  \lean{PhyXMiniProblems.ProblemPhyXMini0810.MassQuantity}
  A nonnegative, unit-independent physical mass
\end{definition}
\begin{proof}
  This is the declared model definition of Mass Quantity.
\end{proof}
\begin{definition}[Vertical Position Quantity]
  \label{def:physics:phyx-mini-0810:phyxminiproblems-problemphyxmini0810-verticalpositionquantity}
  \lean{PhyXMiniProblems.ProblemPhyXMini0810.VerticalPositionQuantity}
  A signed vertical position relative to the coordinate origin in the figure
\end{definition}
\begin{proof}
  This is the declared model definition of Vertical Position Quantity.
\end{proof}
\begin{definition}[Vertical Velocity Quantity]
  \label{def:physics:phyx-mini-0810:phyxminiproblems-problemphyxmini0810-verticalvelocityquantity}
  \lean{PhyXMiniProblems.ProblemPhyXMini0810.VerticalVelocityQuantity}
  \uses{def:physics:phyx-mini-0810:phyxminiproblems-problemphyxmini0810-verticalvelocitydimension}
  A signed component of vertical velocity
\end{definition}
\begin{proof}
  This is the declared model definition of Vertical Velocity Quantity.
\end{proof}
\begin{definition}[Vertical Acceleration Quantity]
  \label{def:physics:phyx-mini-0810:phyxminiproblems-problemphyxmini0810-verticalaccelerationquantity}
  \lean{PhyXMiniProblems.ProblemPhyXMini0810.VerticalAccelerationQuantity}
  \uses{def:physics:phyx-mini-0810:phyxminiproblems-problemphyxmini0810-verticalaccelerationdimension}
  A signed component of vertical acceleration
\end{definition}
\begin{proof}
  This is the declared model definition of Vertical Acceleration Quantity.
\end{proof}
\begin{definition}[Acceleration Magnitude Quantity]
  \label{def:physics:phyx-mini-0810:phyxminiproblems-problemphyxmini0810-accelerationmagnitudequantity}
  \lean{PhyXMiniProblems.ProblemPhyXMini0810.AccelerationMagnitudeQuantity}
  \uses{def:physics:phyx-mini-0810:phyxminiproblems-problemphyxmini0810-verticalaccelerationdimension}
  A nonnegative magnitude of gravitational acceleration
\end{definition}
\begin{proof}
  This is the declared model definition of Acceleration Magnitude Quantity.
\end{proof}
\begin{definition}[mass Readout]
  \label{def:physics:phyx-mini-0810:phyxminiproblems-problemphyxmini0810-massreadout}
  \lean{PhyXMiniProblems.ProblemPhyXMini0810.massReadout}
  \uses{def:physics:phyx-mini-0810:phyxminiproblems-problemphyxmini0810-massquantity}
  Read a physical mass in an arbitrary coherent unit system
\end{definition}
\begin{proof}
  This is the declared model definition of mass Readout.
\end{proof}
\begin{definition}[position Readout]
  \label{def:physics:phyx-mini-0810:phyxminiproblems-problemphyxmini0810-positionreadout}
  \lean{PhyXMiniProblems.ProblemPhyXMini0810.positionReadout}
  \uses{def:physics:phyx-mini-0810:phyxminiproblems-problemphyxmini0810-verticalpositionquantity}
  Read a signed vertical position in an arbitrary coherent unit system
\end{definition}
\begin{proof}
  This is the declared model definition of position Readout.
\end{proof}
\begin{definition}[velocity Readout]
  \label{def:physics:phyx-mini-0810:phyxminiproblems-problemphyxmini0810-velocityreadout}
  \lean{PhyXMiniProblems.ProblemPhyXMini0810.velocityReadout}
  \uses{def:physics:phyx-mini-0810:phyxminiproblems-problemphyxmini0810-verticalvelocityquantity}
  Read a signed vertical velocity in an arbitrary coherent unit system
\end{definition}
\begin{proof}
  This is the declared model definition of velocity Readout.
\end{proof}
\begin{definition}[acceleration Readout]
  \label{def:physics:phyx-mini-0810:phyxminiproblems-problemphyxmini0810-accelerationreadout}
  \lean{PhyXMiniProblems.ProblemPhyXMini0810.accelerationReadout}
  \uses{def:physics:phyx-mini-0810:phyxminiproblems-problemphyxmini0810-verticalaccelerationquantity}
  Read a signed vertical acceleration in an arbitrary coherent unit system
\end{definition}
\begin{proof}
  This is the declared model definition of acceleration Readout.
\end{proof}
\begin{definition}[acceleration Magnitude Readout]
  \label{def:physics:phyx-mini-0810:phyxminiproblems-problemphyxmini0810-accelerationmagnitudereadout}
  \lean{PhyXMiniProblems.ProblemPhyXMini0810.accelerationMagnitudeReadout}
  \uses{def:physics:phyx-mini-0810:phyxminiproblems-problemphyxmini0810-accelerationmagnitudequantity}
  Read a nonnegative acceleration magnitude in coherent units
\end{definition}
\begin{proof}
  This is the declared model definition of acceleration Magnitude Readout.
\end{proof}
\begin{definition}[Motion Event]
  \label{def:physics:phyx-mini-0810:phyxminiproblems-problemphyxmini0810-motionevent}
  \lean{PhyXMiniProblems.ProblemPhyXMini0810.MotionEvent}
  The two endpoint states explicitly drawn in the primary figure
\end{definition}
\begin{proof}
  This is the declared model definition of Motion Event.
\end{proof}
\begin{definition}[Vertical Direction]
  \label{def:physics:phyx-mini-0810:phyxminiproblems-problemphyxmini0810-verticaldirection}
  \lean{PhyXMiniProblems.ProblemPhyXMini0810.VerticalDirection}
  Vertical directions used for the coordinate convention and launch arrow
\end{definition}
\begin{proof}
  This is the declared model definition of Vertical Direction.
\end{proof}
\begin{definition}[Thrown Object Kind]
  \label{def:physics:phyx-mini-0810:phyxminiproblems-problemphyxmini0810-thrownobjectkind}
  \lean{PhyXMiniProblems.ProblemPhyXMini0810.ThrownObjectKind}
  The physical object identified by the problem statement
\end{definition}
\begin{proof}
  This is the declared model definition of Thrown Object Kind.
\end{proof}
\begin{definition}[Vertical Motion Model]
  \label{def:physics:phyx-mini-0810:phyxminiproblems-problemphyxmini0810-verticalmotionmodel}
  \lean{PhyXMiniProblems.ProblemPhyXMini0810.VerticalMotionModel}
  The idealized dynamics selected by the phrase "ignoring air resistance."
\end{definition}
\begin{proof}
  This is the declared model definition of Vertical Motion Model.
\end{proof}
\begin{definition}[Figure Quantity Label]
  \label{def:physics:phyx-mini-0810:phyxminiproblems-problemphyxmini0810-figurequantitylabel}
  \lean{PhyXMiniProblems.ProblemPhyXMini0810.FigureQuantityLabel}
  The five symbolic quantity labels visible in image 810.png
\end{definition}
\begin{proof}
  This is the declared model definition of Figure Quantity Label.
\end{proof}
\begin{definition}[Baseball Vertical Throw Figure]
  \label{def:physics:phyx-mini-0810:phyxminiproblems-problemphyxmini0810-baseballverticalthrowfigure}
  \lean{PhyXMiniProblems.ProblemPhyXMini0810.BaseballVerticalThrowFigure}
  \uses{def:physics:phyx-mini-0810:phyxminiproblems-problemphyxmini0810-motionevent, def:physics:phyx-mini-0810:phyxminiproblems-problemphyxmini0810-verticaldirection, def:physics:phyx-mini-0810:phyxminiproblems-problemphyxmini0810-figurequantitylabel}
  Literal content retained from the primary raster. The numerical fields are displayed SI readouts; y₂ is retained only as a visible symbolic label and is not assigned a numerical value in the figure data
\end{definition}
\begin{proof}
  This is the declared model definition of Baseball Vertical Throw Figure.
\end{proof}
\begin{definition}[Baseball Vertical Throw Setup]
  \label{def:physics:phyx-mini-0810:phyxminiproblems-problemphyxmini0810-baseballverticalthrowsetup}
  \lean{PhyXMiniProblems.ProblemPhyXMini0810.BaseballVerticalThrowSetup}
  \uses{def:physics:phyx-mini-0810:phyxminiproblems-problemphyxmini0810-massquantity, def:physics:phyx-mini-0810:phyxminiproblems-problemphyxmini0810-verticalpositionquantity, def:physics:phyx-mini-0810:phyxminiproblems-problemphyxmini0810-verticalvelocityquantity, def:physics:phyx-mini-0810:phyxminiproblems-problemphyxmini0810-verticalaccelerationquantity, def:physics:phyx-mini-0810:phyxminiproblems-problemphyxmini0810-accelerationmagnitudequantity, def:physics:phyx-mini-0810:phyxminiproblems-problemphyxmini0810-motionevent, def:physics:phyx-mini-0810:phyxminiproblems-problemphyxmini0810-verticaldirection, def:physics:phyx-mini-0810:phyxminiproblems-problemphyxmini0810-thrownobjectkind, def:physics:phyx-mini-0810:phyxminiproblems-problemphyxmini0810-verticalmotionmodel, def:physics:phyx-mini-0810:phyxminiproblems-problemphyxmini0810-baseballverticalthrowfigure}
  Independent physical quantities of the throw. In particular, the position at highestPoint is an unconstrained observable here; it is neither defined from the recorded answer nor from a closed-form height formula
\end{definition}
\begin{proof}
  This is the declared model definition of Baseball Vertical Throw Setup.
\end{proof}
\begin{definition}[mass In Kilograms]
  \label{def:physics:phyx-mini-0810:phyxminiproblems-problemphyxmini0810-massinkilograms}
  \lean{PhyXMiniProblems.ProblemPhyXMini0810.massInKilograms}
  \uses{def:physics:phyx-mini-0810:phyxminiproblems-problemphyxmini0810-massreadout, def:physics:phyx-mini-0810:phyxminiproblems-problemphyxmini0810-baseballverticalthrowsetup}
  Coherent-SI kilogram readout of the baseball mass
\end{definition}
\begin{proof}
  This is the declared model definition of mass In Kilograms.
\end{proof}
\begin{definition}[position In Meters]
  \label{def:physics:phyx-mini-0810:phyxminiproblems-problemphyxmini0810-positioninmeters}
  \lean{PhyXMiniProblems.ProblemPhyXMini0810.positionInMeters}
  \uses{def:physics:phyx-mini-0810:phyxminiproblems-problemphyxmini0810-positionreadout, def:physics:phyx-mini-0810:phyxminiproblems-problemphyxmini0810-motionevent, def:physics:phyx-mini-0810:phyxminiproblems-problemphyxmini0810-baseballverticalthrowsetup}
  Coherent-SI metre readout of the position at a labeled event
\end{definition}
\begin{proof}
  This is the declared model definition of position In Meters.
\end{proof}
\begin{definition}[velocity In Meters Per Second]
  \label{def:physics:phyx-mini-0810:phyxminiproblems-problemphyxmini0810-velocityinmeterspersecond}
  \lean{PhyXMiniProblems.ProblemPhyXMini0810.velocityInMetersPerSecond}
  \uses{def:physics:phyx-mini-0810:phyxminiproblems-problemphyxmini0810-velocityreadout, def:physics:phyx-mini-0810:phyxminiproblems-problemphyxmini0810-motionevent, def:physics:phyx-mini-0810:phyxminiproblems-problemphyxmini0810-baseballverticalthrowsetup}
  Coherent-SI metre-per-second readout at a labeled event
\end{definition}
\begin{proof}
  This is the declared model definition of velocity In Meters Per Second.
\end{proof}
\begin{definition}[vertical Acceleration In Meters Per Second Squared]
  \label{def:physics:phyx-mini-0810:phyxminiproblems-problemphyxmini0810-verticalaccelerationinmeterspersecondsquared}
  \lean{PhyXMiniProblems.ProblemPhyXMini0810.verticalAccelerationInMetersPerSecondSquared}
  \uses{def:physics:phyx-mini-0810:phyxminiproblems-problemphyxmini0810-accelerationreadout, def:physics:phyx-mini-0810:phyxminiproblems-problemphyxmini0810-baseballverticalthrowsetup}
  Coherent-SI metre-per-second-squared readout of signed acceleration
\end{definition}
\begin{proof}
  This is the declared model definition of vertical Acceleration In Meters Per Second Squared.
\end{proof}
\begin{definition}[gravity Magnitude In Meters Per Second Squared]
  \label{def:physics:phyx-mini-0810:phyxminiproblems-problemphyxmini0810-gravitymagnitudeinmeterspersecondsquared}
  \lean{PhyXMiniProblems.ProblemPhyXMini0810.gravityMagnitudeInMetersPerSecondSquared}
  \uses{def:physics:phyx-mini-0810:phyxminiproblems-problemphyxmini0810-accelerationmagnitudereadout, def:physics:phyx-mini-0810:phyxminiproblems-problemphyxmini0810-baseballverticalthrowsetup}
  Coherent-SI metre-per-second-squared readout of the magnitude g
\end{definition}
\begin{proof}
  This is the declared model definition of gravity Magnitude In Meters Per Second Squared.
\end{proof}
\begin{definition}[height Gain In Meters]
  \label{def:physics:phyx-mini-0810:phyxminiproblems-problemphyxmini0810-heightgaininmeters}
  \lean{PhyXMiniProblems.ProblemPhyXMini0810.heightGainInMeters}
  \uses{def:physics:phyx-mini-0810:phyxminiproblems-problemphyxmini0810-baseballverticalthrowsetup, def:physics:phyx-mini-0810:phyxminiproblems-problemphyxmini0810-positioninmeters}
  The requested height is the peak position minus the release position
\end{definition}
\begin{proof}
  This is the declared model definition of height Gain In Meters.
\end{proof}
\begin{definition}[Matches Problem Description]
  \label{def:physics:phyx-mini-0810:phyxminiproblems-problemphyxmini0810-matchesproblemdescription}
  \lean{PhyXMiniProblems.ProblemPhyXMini0810.MatchesProblemDescription}
  \uses{def:physics:phyx-mini-0810:phyxminiproblems-problemphyxmini0810-baseballverticalthrowsetup}
  The prose assumptions: a baseball is thrown upward in ideal uniform gravity
\end{definition}
\begin{proof}
  This is the declared model definition of Matches Problem Description.
\end{proof}
\begin{definition}[Matches Primary Figure]
  \label{def:physics:phyx-mini-0810:phyxminiproblems-problemphyxmini0810-matchesprimaryfigure}
  \lean{PhyXMiniProblems.ProblemPhyXMini0810.MatchesPrimaryFigure}
  \uses{def:physics:phyx-mini-0810:phyxminiproblems-problemphyxmini0810-baseballverticalthrowsetup, def:physics:phyx-mini-0810:phyxminiproblems-problemphyxmini0810-massinkilograms, def:physics:phyx-mini-0810:phyxminiproblems-problemphyxmini0810-positioninmeters, def:physics:phyx-mini-0810:phyxminiproblems-problemphyxmini0810-velocityinmeterspersecond}
  Exact transcription of the primary bitmap. The physical readouts are tied to the displayed numbers, while the unknown label y₂ receives no numerical premise
\end{definition}
\begin{proof}
  This is the declared model definition of Matches Primary Figure.
\end{proof}
\begin{definition}[Uses Standard Earth Gravity]
  \label{def:physics:phyx-mini-0810:phyxminiproblems-problemphyxmini0810-usesstandardearthgravity}
  \lean{PhyXMiniProblems.ProblemPhyXMini0810.UsesStandardEarthGravity}
  \uses{def:physics:phyx-mini-0810:phyxminiproblems-problemphyxmini0810-baseballverticalthrowsetup, def:physics:phyx-mini-0810:phyxminiproblems-problemphyxmini0810-gravitymagnitudeinmeterspersecondsquared}
  The conventional near-Earth value used by the multiple-choice calculation
\end{definition}
\begin{proof}
  This is the declared model definition of Uses Standard Earth Gravity.
\end{proof}
\begin{definition}[Has Physical Throw Parameters]
  \label{def:physics:phyx-mini-0810:phyxminiproblems-problemphyxmini0810-hasphysicalthrowparameters}
  \lean{PhyXMiniProblems.ProblemPhyXMini0810.HasPhysicalThrowParameters}
  \uses{def:physics:phyx-mini-0810:phyxminiproblems-problemphyxmini0810-baseballverticalthrowsetup, def:physics:phyx-mini-0810:phyxminiproblems-problemphyxmini0810-massinkilograms, def:physics:phyx-mini-0810:phyxminiproblems-problemphyxmini0810-velocityinmeterspersecond, def:physics:phyx-mini-0810:phyxminiproblems-problemphyxmini0810-gravitymagnitudeinmeterspersecondsquared}
  Positivity conditions select the physical upward-throw branch without fixing the requested numerical height
\end{definition}
\begin{proof}
  This is the declared model definition of Has Physical Throw Parameters.
\end{proof}
\begin{definition}[Satisfies Uniform Gravity Kinematics]
  \label{def:physics:phyx-mini-0810:phyxminiproblems-problemphyxmini0810-satisfiesuniformgravitykinematics}
  \lean{PhyXMiniProblems.ProblemPhyXMini0810.SatisfiesUniformGravityKinematics}
  \uses{def:physics:phyx-mini-0810:phyxminiproblems-problemphyxmini0810-positionreadout, def:physics:phyx-mini-0810:phyxminiproblems-problemphyxmini0810-velocityreadout, def:physics:phyx-mini-0810:phyxminiproblems-problemphyxmini0810-accelerationreadout, def:physics:phyx-mini-0810:phyxminiproblems-problemphyxmini0810-accelerationmagnitudereadout, def:physics:phyx-mini-0810:phyxminiproblems-problemphyxmini0810-baseballverticalthrowsetup}
  The constant-gravity kinematic laws, stated for every coherent unit choice. They relate the independent endpoint observables and contain neither 1000 / 49, the rounded value 20.4, nor any answer-choice label
\end{definition}
\begin{proof}
  This is the declared model definition of Satisfies Uniform Gravity Kinematics.
\end{proof}
\begin{definition}[Answer Choice]
  \label{def:physics:phyx-mini-0810:phyxminiproblems-problemphyxmini0810-answerchoice}
  \lean{PhyXMiniProblems.ProblemPhyXMini0810.AnswerChoice}
  Labels of the four displayed multiple-choice answers
\end{definition}
\begin{proof}
  This is the declared model definition of Answer Choice.
\end{proof}
\begin{definition}[displayed Height In Meters]
  \label{def:physics:phyx-mini-0810:phyxminiproblems-problemphyxmini0810-answerchoice-displayedheightinmeters}
  \lean{PhyXMiniProblems.ProblemPhyXMini0810.AnswerChoice.displayedHeightInMeters}
  \uses{def:physics:phyx-mini-0810:phyxminiproblems-problemphyxmini0810-answerchoice}
  Height in metres printed beside each answer choice
\end{definition}
\begin{proof}
  This is the declared model definition of displayed Height In Meters.
\end{proof}
\begin{definition}[recorded Answer Choice]
  \label{def:physics:phyx-mini-0810:phyxminiproblems-problemphyxmini0810-recordedanswerchoice}
  \lean{PhyXMiniProblems.ProblemPhyXMini0810.recordedAnswerChoice}
  \uses{def:physics:phyx-mini-0810:phyxminiproblems-problemphyxmini0810-answerchoice}
  The answer label recorded in the supplied dataset metadata
\end{definition}
\begin{proof}
  This is the declared model definition of recorded Answer Choice.
\end{proof}
\begin{definition}[Fits Displayed Height]
  \label{def:physics:phyx-mini-0810:phyxminiproblems-problemphyxmini0810-fitsdisplayedheight}
  \lean{PhyXMiniProblems.ProblemPhyXMini0810.FitsDisplayedHeight}
  \uses{def:physics:phyx-mini-0810:phyxminiproblems-problemphyxmini0810-baseballverticalthrowsetup, def:physics:phyx-mini-0810:phyxminiproblems-problemphyxmini0810-heightgaininmeters, def:physics:phyx-mini-0810:phyxminiproblems-problemphyxmini0810-answerchoice}
  A tenth-metre display represents the ideal height when the discrepancy is strictly less than half of the last displayed metre digit
\end{definition}
\begin{proof}
  This is the declared model definition of Fits Displayed Height.
\end{proof}
\begin{lemma}[height Gain exact]
  \label{lem:physics:phyx-mini-0810:phyxminiproblems-problemphyxmini0810-heightgain-exact}
  \lean{PhyXMiniProblems.ProblemPhyXMini0810.heightGain_exact}
  \uses{def:physics:phyx-mini-0810:phyxminiproblems-problemphyxmini0810-baseballverticalthrowsetup, def:physics:phyx-mini-0810:phyxminiproblems-problemphyxmini0810-heightgaininmeters, def:physics:phyx-mini-0810:phyxminiproblems-problemphyxmini0810-matchesproblemdescription, def:physics:phyx-mini-0810:phyxminiproblems-problemphyxmini0810-matchesprimaryfigure, def:physics:phyx-mini-0810:phyxminiproblems-problemphyxmini0810-usesstandardearthgravity, def:physics:phyx-mini-0810:phyxminiproblems-problemphyxmini0810-hasphysicalthrowparameters, def:physics:phyx-mini-0810:phyxminiproblems-problemphyxmini0810-satisfiesuniformgravitykinematics}
  Under the ideal model, the exact height gain is 20² / (2 * 9.8)
\end{lemma}
\begin{proof}
  The conclusion follows from the governing relations and figure readouts named in the statement of height Gain exact.
\end{proof}
% --- Archon declaration topology end ---
% --- Archon physics formalization source end ---
